\section*{الفصل \ref{ch:g6:solutions-and-solubility} --- المحاليل والذوبانية}

\begin{solution}{exo:g6:solutions-and-solubility:1}
\omterm{def:g6:solutions-and-solubility:solute}{المذاب}: السكر. \omterm{def:g6:solutions-and-solubility:solute}{المذيب}: الماء. نعم، \omterm{def:g6:solutions-and-solubility:solute}{فالمذيب} هو الماء: إنه \omterm{def:g6:solutions-and-solubility:solute}{محلول
مائي}.
\end{solution}

\begin{solution}{exo:g6:solutions-and-solubility:2}
\omterm{def:g2:mixing-and-dissolving:dissolve}{محلول} لا يستطيع أن يذيب مزيدًا من مذابه. \omterm{def:g6:solutions-and-solubility:solute}{فالمذاب} الذي يُضاف إليه يبقى
غير ذائب في القاع.
\end{solution}

\begin{solution}{exo:g6:solutions-and-solubility:3}
$250 + 20 = \qty{270}{g}$.
\end{solution}

\begin{solution}{exo:g6:solutions-and-solubility:4}
الإيثانول: \omterm{def:g6:solutions-and-solubility:miscible}{قابل للامتزاج}. زيت الطهي: \omterm{def:g6:solutions-and-solubility:miscible}{غير قابل للامتزاج}. الخل: \omterm{def:g6:solutions-and-solubility:miscible}{قابل
للامتزاج} (فهو في معظمه ماء).
\end{solution}

\begin{solution}{exo:g6:solutions-and-solubility:5}
البروبانون (الأسيتون). GHS02: شديد القابلية للاشتعال؛ GHS07: يهيّج العينين،
وقد يسبب بخاره النعاس أو الدوار.
\end{solution}

\begin{solution}{exo:g6:solutions-and-solubility:6}
$360 \times \frac{500}{1000} = \qty{180}{g}$.
\end{solution}

\begin{solution}{exo:g6:solutions-and-solubility:7}
\qty{200}{g} من الماء تذيب على الأكثر $360 \times \frac{200}{1000} =
\qty{72}{g}$. لا: يبقى $80 - 72 = \qty{8}{g}$ في القاع.
\end{solution}

\begin{solution}{exo:g6:solutions-and-solubility:8}
كتلة \omterm{def:g2:mixing-and-dissolving:dissolve}{المحلول} $225 + 25 = \qty{250}{g}$؛ و$\frac{25}{250} =
\qty{10}{\percent}$.
\end{solution}

\begin{solution}{exo:g6:solutions-and-solubility:9}
في الكأس الأولى، حيث ذاب الملح كله ولم يصر \omterm{def:g2:mixing-and-dissolving:dissolve}{المحلول} مشبعًا
بعد.
\end{solution}

\begin{solution}{exo:g6:solutions-and-solubility:10}
كان ثنائي أكسيد الكربون ذائبًا تحت الضغط في القارورة المغلقة. وفتحها
يخفض الضغط: فيصير الماء قادرًا على حمل غاز أقل، ويخرج الغاز الزائد
فقاعاتٍ.
\end{solution}

\begin{solution}{exo:g6:solutions-and-solubility:11}
$2000 \div 0.04 = 50\,000$ مرة.
\end{solution}

\begin{solution}{exo:g6:solutions-and-solubility:12}
الغاز \omterm{def:g2:mixing-and-dissolving:dissolve}{يذوب} في الماء الدافئ أقل: فالبركة الدافئة تحوي أكسجينًا ذائبًا أقل،
فتعاني الأسماك نقص الأكسجين.
\end{solution}

\begin{solution}{pb:g6:solutions-and-solubility:1}
\textbf{1.} \omterm{def:g6:solutions-and-solubility:solute}{المذاب}: الملح. \omterm{def:g6:solutions-and-solubility:solute}{المذيب}: الماء.

\textbf{2.} يحمل الماء أكبر قدر ممكن من الملح: فلا يستطيع أن يذيب فيه
مزيدًا من الملح.

\textbf{3.} الملح غير الذائب في القاع دليل على أن \omterm{def:g2:mixing-and-dissolving:dissolve}{المحلول} الملحي مشبع:
فهو يحمل أكبر قدر ممكن من الملح.

\textbf{4.} $360 \times 20 = \qty{7200}{g} = \qty{7.2}{kg}$.

\textbf{5.} $20 + 7.2 = \qty{27.2}{kg}$.

\textbf{6.} $\frac{7.2}{27.2} \approx 0.26$: أي نحو \qty{26}{\percent}.

\textbf{7.} لا: فالماء يستطيع أن يذيب \qty{7.2}{kg}. ولا يزال يستطيع أن
يذيب $7.2 - 5 = \qty{2.2}{kg}$.

\textbf{8.} $20 - 5 = \qty{15}{kg}$ من الماء (\qty{15}{L}).

\textbf{9.} $360 \times 15 = \qty{5400}{g} = \qty{5.4}{kg}$.

\textbf{10.} يخرج من \omterm{def:g2:mixing-and-dissolving:dissolve}{المحلول} بلوراتٍ صلبة: يتبلور ويهبط إلى
القاع.

\textbf{11.} تظهر بلورات بيضاء جديدة وتتراكم في قاع
الخزان.

\textbf{12.} $7.2 - 5.4 = \qty{1.8}{kg}$ من الملح تبلورت.
\end{solution}
